\section{A resource-transition reading of the datum}\label{sec:benchmark}

The witness datum is abstract by design. To show that the theorems bear on a textbook dynamical system, and to let each symbol be read as a quantity an agency monitors, we realise the datum as a \citet{Schaefer1954} biomass--yield model. The realisation is an exact rational re-reading, not a fitted case study: every number is re-derived by the deposited benchmark script, and no empirical claim is made (full detail in SI~S5).

\paragraph{The system} Let $B(t)$ be spawning biomass (kt) over one annual review period, with logistic surplus production $\varsigma(B)=rB(1-B/K)$, $r=4$, $K=10$, and a biomass limit $B_{\mathrm{lim}}=2$ kt below which recruitment is judged impaired. The ecological floor is the margin $s_1=B-B_{\mathrm{lim}}\ge0$; the income floor $s_2\ge0$ is the fleet income margin net of costs, maintained by the quota and by a transition reserve $x$ (an adjustment fund financing fleet draw-down). The four plans are standard in transition management: \act{no-switch} (status-quo harvest), \act{fast} (an immediate large quota cut -- a pulse season followed by a closure -- accepting a mid-season trough on the ecological floor), \act{slow} (a gradual reduction shifting the burden to the income floor) and \act{staged} (a below-sustained-yield quota while the reserve finances the fleet transition, after which both margins improve). The disturbance is a mid-season heatwave strike imposing an excess mortality $\delta_0=1/2$ kt, present or absent. The state is observed through the institutional dashboard $w\cdot(s_1,s_2)$, and the reviewer checks the dashboard rather than the floors -- a deliberately adverse reporting architecture adopted to isolate the operator issue, not a description of current fisheries practice, in which a limit-reference-point breach triggers a mandatory response.

\paragraph{The schedule, exactly} At the witness point the stock is $B=16/5$ kt. \act{fast} is realised by the quota $H^*(t)=3+\varsigma(B_{\mathrm{line}}(t))$ on $(0,\tfrac12)$ -- a pulse declining from $1463/125\approx11.7$ to $2161/250\approx8.6$ kt/yr, admissible against a fleet capacity of $13$ -- followed by a closure, $H^*=0$. The benign branch tracks the tube line $s_1:\ 6/5\to-3/10\to6/5$ exactly; the heatwave drops the trough from $17/10$ to $6/5$ kt, the adverse tube's value, and on the recovery leg the surplus is at least $\varsigma(6/5)=528/125>4$ kt/yr while the certified recovery needs only $4$, so the realised nonlinear trajectory lies on or above the certified tube throughout: the certified piecewise-linear tubes are conservative for the Schaefer realisation, by an exact inequality that uses the monotonicity of $\varsigma$ below $K/2$. \act{no-switch} and \act{slow} hold $B=16/5$ at the sustained yield $\varsigma(16/5)=1088/125\approx8.7$ kt/yr (below the model's maximum sustainable yield $rK/4=10$); \act{staged}, entered at the rescue witness with fund $x=3/2$, rebuilds the stock $16/5\to69/20$ while the fund draws down $3/2\to1/2$ and both margins improve by $e=1/4$.

\paragraph{The dashboard reading} Under the heatwave branch of the licensed plan the equal-weight index $s_1+s_2$ takes the tube values $(12/5,\ 2/5,\ 12/5)$: the dashboard never drops below $2/5$, and the transition is certified at every weighting (the thresholds $\rho_1=2/3$ and $\rho_2=3/2$ place \act{fast} in every income-heavy assessor's acceptance set and \act{slow} in every biomass-heavy assessor's). The biomass floor, meanwhile, is breached mid-transition on both branches ($-4/5$ with the strike, $-3/10$ without): the dashboard certifies a transition that violates the biomass limit whether or not the heatwave occurs, and no weighting can detect this from the index alone (Figure~\ref{fig:benchmark}). The operational number is the reserve threshold: a fund below the buy-back cost $c=1$ makes the staged transition unfinanceable, so every state of the failure set carries the shortfall $\kappa^*=1-x$, and a top-up of the adjustment fund converts an index-licensed impossibility into a rescue. In management language, the dashboard is structurally unable to distinguish a transition that survives the heatwave from one that does not; the financing gap, not the weighting, is what separates them.

\paragraph{Decision-analytic structure} Read as a decision problem the benchmark is a sorting problem in multi-criteria decision aiding: the decision maker is the resource authority; the alternatives are the four management schedules; the criteria are the two floors, evaluated along the whole transition under the declared disturbance; and the weights are the preference orderings of the doctrine under which the index is read. The criteria conflict -- each principal schedule protects one at the expense of the other -- and each is defended by a distinct constituency, so the weight space reads both as an assessor's calibration and as the stakeholder dispute. The deliverables are ones a decision analyst can communicate: the licensing thresholds (the preference range over which each alternative is certified), the rescue threshold (the adjustment funding that converts an unlicensable state into a staged transition), and, from Theorem~\ref{thm:ladder}, the elasticity below which an aggregator exposes the breach ($\sigma^*\in(2,3)$ at the witness point: every $\sigma\ge3$ aggregator certifies with no mid-transition response, every $\sigma\le2$ member rejects and triggers the closure-and-rebuilding response). The aggregator choice -- a reporting decision -- flips a management action.

\paragraph{Units on a real stock} The units are anchored to a regulated floor: for Northern cod (NAFO Divisions 2J3KL) the assessed spawning-stock series carries a limit reference point of $884.6$ kt \citep{DFO2016}, and reading the benchmark's floor as this LRP places the witness's opening stock ($1.6$ LRP units, $\approx1.42$ Mt) in the neighbourhood of the assessed mid-1980s biomass and the adverse trough ($0.6$ LRP units, $\approx531$ kt) below the LRP but above the 2015 assessed level of about a third of it; the adverse plan's worst-case excursion, one full LRP unit, is of the order of the stock's historical swing across the LRP. The anchoring fixes units only: no parameter of the datum is estimated from data, and the companion scored forecast test on the same series \citep{Abaee2026cod} is a caution of exactly this paper's discipline -- no structural surplus-production model in its ladder beat last-value persistence, so what survives scrutiny on such a stock is the menu and the floors, both set by management and both taken as given here.

\begin{figure}[!htbp]
\centering
\includegraphics[width=0.96\linewidth]{fig_benchmark.png}
\caption{The witness datum realised as a fishery transition. (a) Under the pulse-and-closure plan the composite index $s_1+s_2$ at $w=(1,1)$ remains certified at every instant (minimum $2/5$) while the biomass margin breaches its floor mid-transition ($-4/5$ with the heatwave strike, $-3/10$ without). (b) The plan menu as management schedules: pulse plus closed season; sustained-yield status quo; and the reserve-financed staged transition, which rebuilds the stock while improving both margins. All values are the exact rationals of SI~S5.}
\label{fig:benchmark}
\end{figure}

